\section{Calculus and Analysis}

The differential $\dd$ is upright; $\varepsilon$ is the default small positive quantity in analysis. Function spaces take their domains in parentheses, for example $C([a,b])$, $\mathcal{R}([a,b])$, and $L^p(\Omega)$.

\begin{notationtable}
$\displaystyle\lim_{x\to a} f(x)$ & Limit as $x$ approaches $a$. \\
$\displaystyle\lim_{x\to a^-} f(x)$ & Left-hand limit as $x$ approaches $a$. \\*
$\displaystyle\lim_{x\to a^+} f(x)$ & Right-hand limit as $x$ approaches $a$. \\
$\displaystyle\lim_{x\to\infty} f(x),\ \lim_{x\to-\infty} f(x)$ & Limits as $x$ tends to positive or negative infinity. \\
$\begin{aligned} &\displaystyle\lim_{x\to a} f(x)=\infty,\\[-0.15em] &\displaystyle\lim_{x\to a} f(x)=-\infty \end{aligned}$ & Infinite limits. \\
$\min A$ & Minimum of $A$. \\*
$\max A$ & Maximum of $A$. \\
$\inf A$ & Infimum of $A$. \\*
$\sup A$ & Supremum of $A$. \\
$\overline{\R}=\R\cup\{-\infty,+\infty\}$ & Extended real line. \\
\addlinespace[0.35em]
$f'(a)$ & First derivative of $f$ at $a$. \\
$f^{(n)}(a)$ & $n$th derivative of $f$ at $a$. \\
$\displaystyle\frac{\dd y}{\dd x}$ & Derivative of $y$ with respect to $x$. \\
$\displaystyle\frac{\dd^n y}{\dd x^n}$ & $n$th derivative of $y$ with respect to $x$ in Leibniz notation. \\
$\displaystyle\frac{\partial f}{\partial x_i},\ \partial_i f,\ f_{x_i}$ & Partial derivative with respect to $x_i$. \\
$\displaystyle\frac{\partial^2 f}{\partial x_i\partial x_j},\ \partial_{ij}f,\ f_{x_i x_j}$ & Second or mixed partial derivative. \\
$D_{\mathbf v}f(\mathbf a),\ \partial_{\mathbf v}f(\mathbf a)$ & Directional derivative along the vector $\mathbf v$, which need not be a unit vector: $\displaystyle D_{\mathbf v}f(\mathbf a)=\left.\frac{\dd}{\dd t}f(\mathbf a+t\mathbf v)\right|_{t=0}=Df(\mathbf a)[\mathbf v]=\nabla f(\mathbf a)\cdot\mathbf v$ when $f$ is differentiable. For the unit direction determined by $\mathbf v\ne\mathbf0$, use $\mathbf u=\mathbf v/\|\mathbf v\|$. \\
$\displaystyle\dd f=\sum_{i=1}^n\frac{\partial f}{\partial x_i}\,\dd x_i$ & Total differential of a scalar function of $n$ variables. \\
\addlinespace[0.35em]
$\displaystyle\int f(x)\,\dd x$ & Indefinite integral; a family of antiderivatives, with an additive constant understood. \\
$\displaystyle\int_a^b f(x)\,\dd x$ & Definite integral over $[a,b]$. \\
$\begin{aligned} &\displaystyle\int_a^\infty f(x)\,\dd x,\\[-0.15em] &\displaystyle\int_{-\infty}^b f(x)\,\dd x,\ \int_{-\infty}^{\infty}f(x)\,\dd x \end{aligned}$ & Improper integrals over unbounded intervals. \\
$\displaystyle\int_a^b f\,\dd\alpha$ & Riemann--Stieltjes integral of $f$ with respect to $\alpha$. \\
$\begin{aligned} &P=\{x_0,\dots,x_n\},\\[-0.15em] &a=x_0<\cdots<x_n=b \end{aligned}$ & Partition of the interval $[a,b]$. \\
$\begin{aligned} &\displaystyle\sum_{i=1}^n f(\xi_i)\Delta x_i,\\[-0.15em] &\Delta x_i=x_i-x_{i-1} \end{aligned}$ & Tagged Riemann sum, with $\xi_i\in[x_{i-1},x_i]$. \\
$\displaystyle\iint_D f\,\dd A$ & Double integral over $D$. \\
$\displaystyle\iiint_\Omega f\,\dd V$ & Triple integral over $\Omega$. \\
$\displaystyle\int_C f(z)\,\dd z$ & Contour integral along an oriented contour $C$. \\
$\displaystyle\oint_C f(z)\,\dd z$ & Contour integral over a closed contour $C$. \\
$-C$ & The contour $C$ with the opposite orientation. \\
$\gamma^{-1}$ & Reversed parametrized path $\gamma$, when unambiguous; not a functional inverse. \\
$\PV,\ \mathrm{P.V.}$ & Principal value; in particular, the Cauchy principal value of a singular or improper integral. \\
\addlinespace[0.35em]
$\displaystyle\sum_{k=1}^{n} a_k$ & Finite sum. \\
$\displaystyle\sum_{k=1}^{\infty} a_k$ & Infinite series. \\
$\displaystyle\prod_{k=1}^{n} a_k$ & Finite product. \\
$\displaystyle\sum_{n=0}^{\infty}a_n(x-a)^n$ & Power series centered at $a$. \\
$\displaystyle\sum_{k=0}^{n}\frac{f^{(k)}(a)}{k!}(x-a)^k$ & $n$th Taylor polynomial of $f$ about $a$; no single letter for the polynomial is imposed globally. \\
\addlinespace[0.35em]
$e^x,\ \exp x$ & Exponential function. \\
$\ln x,\ \log_e x$ & Logarithm to base $e$. \\
$\lg x,\ \log_{10} x$ & Logarithm to base $10$. \\
$\lb x,\ \log_2 x$ & Logarithm to base $2$. \\
$\log_a x$ & Logarithm to base $a$. \\
\addlinespace[0.35em]
$\sin x$ & Sine. \\*
$\cos x$ & Cosine. \\*
$\tan x$ & Tangent. \\*
$\csc x$ & Cosecant. \\*
$\sec x$ & Secant. \\*
$\cot x$ & Cotangent. \\
\addlinespace[0.35em]
$\arcsin x,\ \sin^{-1}x$ & Principal inverse sine, $[-1,1]\to[-\pi/2,\pi/2]$. \\*
$\arccos x,\ \cos^{-1}x$ & Principal inverse cosine, $[-1,1]\to[0,\pi]$. \\*
$\arctan x,\ \tan^{-1}x$ & Principal inverse tangent, $\R\to(-\pi/2,\pi/2)$. \\*
$\arccsc x,\ \csc^{-1}x$ & Principal inverse cosecant, defined by $\arccsc x=\arcsin(1/x)$ for $|x|\ge1$; its range is $[-\pi/2,0)\cup(0,\pi/2]$. \\*
$\arcsec x,\ \sec^{-1}x$ & Principal inverse secant, defined by $\arcsec x=\arccos(1/x)$ for $|x|\ge1$; its range is $[0,\pi/2)\cup(\pi/2,\pi]$. \\*
$\arccot x,\ \cot^{-1}x$ & Principal real inverse cotangent, $\R\to(0,\pi)$. \\
\addlinespace[0.35em]
$\sinh x$ & Hyperbolic sine. \\*
$\cosh x$ & Hyperbolic cosine. \\*
$\tanh x$ & Hyperbolic tangent. \\*
$\csch x$ & Hyperbolic cosecant. \\*
$\sech x$ & Hyperbolic secant. \\*
$\coth x$ & Hyperbolic cotangent. \\
\addlinespace[0.35em]
$\arcsinh x,\ \sinh^{-1}x$ & Principal real inverse hyperbolic sine, $\R\to\R$. \\*
$\arccosh x,\ \cosh^{-1}x$ & Principal real inverse hyperbolic cosine, $[1,\infty)\to[0,\infty)$. \\*
$\arctanh x,\ \tanh^{-1}x$ & Principal real inverse hyperbolic tangent, $(-1,1)\to\R$. \\*
$\arccsch x,\ \csch^{-1}x$ & Principal real inverse hyperbolic cosecant, defined by $\arccsch x=\arcsinh(1/x)$ on $\R\setminus\{0\}$; its range is $\R\setminus\{0\}$. \\*
$\arcsech x,\ \sech^{-1}x$ & Principal real inverse hyperbolic secant, defined by $\arcsech x=\arccosh(1/x)$ on $(0,1]$; its range is $[0,\infty)$. \\*
$\arccoth x,\ \coth^{-1}x$ & Principal real inverse hyperbolic cotangent, defined by $\arccoth x=\arctanh(1/x)$ for $|x|>1$; its range is $\R\setminus\{0\}$. \\
\addlinespace[0.35em]
$\Gamma(z)$ & Gamma function. \\
$B(x,y)$ & Beta function. \\
$\zeta(s)$ & Riemann zeta function. \\
\addlinespace[0.35em]
$\mathbf J_{\mathbf F}(\mathbf a),\ D\mathbf F(\mathbf a)$ & Jacobian matrix of $\mathbf F=(F_1,\dots,F_m)\colon\R^n\to\R^m$ at $\mathbf a$, with $\displaystyle\mathbf J_{\mathbf F}(\mathbf a)=\left(\frac{\partial F_i}{\partial x_j}(\mathbf a)\right)_{i,j}$. The notation $D\mathbf F(\mathbf a)$ denotes the total derivative, identified with this matrix in standard coordinates. \\
$\begin{aligned} &\det\mathbf J_{\mathbf F}(\mathbf a),\ \det D\mathbf F(\mathbf a),\\[-0.15em] &\displaystyle\left.\frac{\partial(F_1,\dots,F_n)}{\partial(x_1,\dots,x_n)}\right|_{\mathbf a} \end{aligned}$ & Jacobian determinant of $\mathbf F\colon\R^n\to\R^n$ at $\mathbf a$, in preferred order. \\
$\Hess f$ & Hessian matrix of $f$. \\
$\nabla f,\ \operatorname{grad}f$ & Gradient of $f$. \\
$\nabla\cdot\mathbf{F},\ \operatorname{div}\mathbf F$ & Divergence of $\mathbf{F}$. \\
$\nabla\times\mathbf{F},\ \operatorname{curl}\mathbf F$ & Curl of $\mathbf{F}$. \\
$\Delta f$ & Laplacian of $f$. \\
$\mathbf{r}(t)$ & Parametrized curve. \\
$\dd s$ & Element of arc length. \\
$\displaystyle\int_C f\,\dd s$ & Scalar line integral along $C$ with respect to arc length. \\
$\displaystyle\int_C\mathbf{F}\cdot\dd\mathbf{r}$ & Line integral of the vector field $\mathbf F$ along $C$. \\
$\displaystyle\iint_S f\,\dd S$ & Scalar surface integral over $S$. \\
$\displaystyle\iint_S\mathbf{F}\cdot\mathbf{n}\,\dd S$ & Flux integral through $S$. \\
\addlinespace[0.35em]
$x_n\to x,\ \displaystyle\lim_{n\to\infty}x_n=x$ & Convergence of a sequence. \\
$(x_{n_k})$ & Subsequence of the sequence $(x_n)$. \\
$a_n\nearrow L$ & Monotone increasing convergence to $L$. \\*
$a_n\searrow L$ & Monotone decreasing convergence to $L$. \\
$\limsup a_n$ & Limit superior. \\*
$\liminf a_n$ & Limit inferior. \\
$f(x)\sim g(x)$ & Asymptotic equivalence in the stated limiting regime, typically meaning $f(x)/g(x)\to1$. \\
$f_n\ptwto f,\ f_n\ptwlabelto f$ & Pointwise convergence of functions on a common domain by default when no ambient topology, metric, or norm is specified. If such a structure is specified, the ordinary arrow $\to$ refers to convergence in that structure. \\*
$f_n\unifto f,\ f_n\uniflabelto f$ & Uniform convergence. \\
$f_n\aeto f$ & Convergence almost everywhere. \\
\addlinespace[0.35em]
$\mu(E)$ & Measure of a measurable set $E$. \\
$\displaystyle\int_E f\,\dd\mu$ & Integral of $f$ over $E$ with respect to the measure $\mu$. \\
$f=g\ \text{a.e.},\quad f=g\ \mu\text{-a.e.}$ & Equality almost everywhere, with the measure omitted when understood. \\
$\displaystyle\esssup_{x\in E} f(x)$ & Essential supremum of $f$ on $E$. \\
$\displaystyle\essinf_{x\in E} f(x)$ & Essential infimum of $f$ on $E$. \\
\addlinespace[0.35em]
$C(\Omega),\ C^0(\Omega)$ & Continuous functions on $\Omega$. \\*
$C^k(\Omega)$ & $k$ times continuously differentiable functions on $\Omega$. \\*
$C^\infty(\Omega)$ & Smooth functions on $\Omega$. \\
$\operatorname{supp}f$ & Support of $f$, the closure of the set on which $f$ is nonzero. \\
$C_c(\Omega),\ C_c^\infty(\Omega)$ & Continuous and smooth functions on $\Omega$ with compact support, respectively. \\
$\mathcal{R}([a,b])$ & Riemann-integrable functions on $[a,b]$. \\
$L^p(\Omega),\ L^p(X,\mu)$ & Lebesgue $L^p$ spaces; the measure is displayed when it is not understood. \\*
$L^\infty(\Omega),\ L^\infty(X,\mu)$ & Lebesgue spaces of essentially bounded functions. \\
$\ell^p$ & Space of scalar sequences $(x_n)$ with $\sum_{n=1}^{\infty}|x_n|^p<\infty$, for $1\le p<\infty$. \\
$\ell^\infty$ & Space of bounded scalar sequences. \\
$\|f\|_\infty$ & Supremum norm $\displaystyle\sup_{x\in\Omega}|f(x)|$ for an actual bounded function, especially in bounded- or continuous-function spaces. \\
$\|f\|_{L^p(\Omega)},\ \|f\|_p$ & $L^p$-norm $\displaystyle\left(\int_\Omega |f(x)|^p\,\dd x\right)^{1/p}$ for $1\le p<\infty$. \\
$\|f\|_{L^\infty(\Omega)}$ & $L^\infty$-norm $\displaystyle\esssup_{x\in\Omega}|f(x)|$; this notation is preferred when distinguishing the essential-supremum norm from an ordinary supremum norm. \\
\addlinespace[0.35em]
$\alpha=(\alpha_1,\dots,\alpha_n)\in\Nzero^n$ & Multi-index; $\alpha$ is not treated as a vector for typographic purposes. \\
$\begin{aligned} &|\alpha|=\alpha_1+\cdots+\alpha_n,\\[-0.15em] &\alpha!=\alpha_1!\cdots\alpha_n! \end{aligned}$ & Order and factorial of a multi-index. \\
$\mathbf x^\alpha=x_1^{\alpha_1}\cdots x_n^{\alpha_n}$ & Multi-index monomial for $\mathbf x=(x_1,\dots,x_n)$. \\
$D^\alpha f,\ \partial^\alpha f$ & Mixed partial derivative specified by the multi-index $\alpha$. \\
\addlinespace[0.35em]
$B_r(x)$ & Open ball of radius $r$ centered at $x$. \\*
$\overline{B}_r(x)$ & Closed ball of radius $r$ centered at $x$. \\
\addlinespace[0.35em]
$\widehat{\C}=\C\cup\{\infty\}$ & Extended complex plane, equivalently the Riemann sphere. \\
$z=x+iy$ & Complex number with real part $x$ and imaginary part $y$. \\
$\RePart z,\ \Re z$ & Real part of $z$. \\
$\ImPart z,\ \Im z$ & Imaginary part of $z$; distinct from lowercase $\im$ for image. \\
$\overline{z}$ & Complex conjugate. \\
$|z|$ & Complex modulus. \\
$\arg z$ & Argument of a nonzero complex number $z$; its values differ by integer multiples of $2\pi$. \\
$\Arg z$ & Principal argument of $z\ne0$, with $-\pi<\Arg z\le\pi$. \\
$\Ln z$ & General multivalued complex logarithm for $z\ne0$: $\displaystyle\Ln z=\{\ln|z|+i(\Arg z+2\pi k)\mid k\in\Z\}$. \\
$\ln z=\ln|z|+i\Arg z$ & Principal value of the complex logarithm for $z\ne0$. On $\C\setminus(-\infty,0]$, it is the analytic principal branch, with $-\pi<\Arg z<\pi$. \\
$\Res(f;a),\ \operatorname{Res}_{z=a}f(z)$ & Residue of $f$ at a finite point $a$. \\
$\Res(f;\infty),\ \operatorname{Res}_{z=\infty}f(z)$ & Residue of $f$ at infinity. \\
\end{notationtable}
